% paper45-callable-surfaces.tex — Callable Surface Analysis
% Paper 45 of the Judgment Geometry series.
% Compile: pdflatex -interaction=nonstopmode paper45-callable-surfaces.tex
\documentclass[11pt]{article}
\usepackage{lmodern}
% jugeo-common.tex — shared preamble for all Judgment Geometry papers
% Include via: % jugeo-common.tex — shared preamble for all Judgment Geometry papers
% Include via: % jugeo-common.tex — shared preamble for all Judgment Geometry papers
% Include via: \input{jugeo-common}

% ── Packages ────────────────────────────────────────────────────────────
\usepackage{amsmath,amssymb,amsthm,mathtools}
\usepackage{tikz-cd}
\usepackage{listings}
\usepackage{xcolor}
\usepackage{xspace}
\usepackage{booktabs}
\usepackage{multirow}
\usepackage{enumitem}
\usepackage[expansion=false]{microtype}
\usepackage{hyperref}
\usepackage{cleveref}
% \usepackage{thmtools}  % disabled: not available in this TeX installation
\usepackage{float}
\usepackage{caption}
\usepackage{subcaption}
\usepackage{tabularx}
\usepackage{stmaryrd}       % \llbracket, \rrbracket
\usepackage{bbm}            % \mathbbm{1}

% ── Page geometry ───────────────────────────────────────────────────────
\usepackage[margin=1in]{geometry}

% ── Theorem environments ────────────────────────────────────────────────
\theoremstyle{plain}
\newtheorem{theorem}{Theorem}[section]
\newtheorem{lemma}[theorem]{Lemma}
\newtheorem{proposition}[theorem]{Proposition}
\newtheorem{corollary}[theorem]{Corollary}

\theoremstyle{definition}
\newtheorem{definition}[theorem]{Definition}
\newtheorem{example}[theorem]{Example}
\newtheorem{construction}[theorem]{Construction}

\theoremstyle{remark}
\newtheorem{remark}[theorem]{Remark}
\newtheorem{notation}[theorem]{Notation}

% ── Python listings style ──────────────────────────────────────────────
\definecolor{jugeoBlue}{HTML}{2563EB}
\definecolor{jugeoGreen}{HTML}{059669}
\definecolor{jugeoRed}{HTML}{DC2626}
\definecolor{jugeoPurple}{HTML}{7C3AED}
\definecolor{jugeoGray}{HTML}{6B7280}
\definecolor{jugeoBg}{HTML}{F8FAFC}

\lstdefinestyle{jugeo-python}{
  language=Python,
  basicstyle=\ttfamily\small,
  keywordstyle=\color{jugeoBlue}\bfseries,
  stringstyle=\color{jugeoGreen},
  commentstyle=\color{jugeoGray}\itshape,
  numberstyle=\tiny\color{jugeoGray},
  backgroundcolor=\color{jugeoBg},
  numbers=left,
  numbersep=6pt,
  frame=single,
  framerule=0.4pt,
  rulecolor=\color{jugeoGray!40},
  breaklines=true,
  breakatwhitespace=true,
  showstringspaces=false,
  tabsize=4,
  xleftmargin=1.5em,
  framexleftmargin=1.5em,
  morekeywords={dataclass,frozen,field,Enum,IntEnum,auto,Any,
                Mapping,Sequence,Optional,match,case},
  literate={→}{{$\to$}}1 {←}{{$\leftarrow$}}1
           {≤}{{$\leq$}}1 {≥}{{$\geq$}}1 {≠}{{$\neq$}}1
           {∈}{{$\in$}}1 {∀}{{$\forall$}}1 {∃}{{$\exists$}}1
           {⊕}{{$\oplus$}}1 {⊖}{{$\ominus$}}1,
  escapeinside={(*@}{@*)},
}
\lstset{style=jugeo-python}

\lstdefinestyle{jugeo-output}{
  basicstyle=\ttfamily\small,
  backgroundcolor=\color{jugeoBg},
  frame=single,
  framerule=0.4pt,
  rulecolor=\color{jugeoGray!40},
  breaklines=true,
  showstringspaces=false,
  xleftmargin=1.5em,
  framexleftmargin=0.5em,
  numbers=none,
}

% ── JuGeo-specific macros ──────────────────────────────────────────────

% Judgment 8-tuple
\newcommand{\J}{\mathcal{J}}
\newcommand{\Jud}[8]{(#1,\,#2,\,#3,\,#4,\,#5,\,#6,\,#7,\,#8)}
\newcommand{\JudShort}[1]{\J_{#1}}

% Components
\newcommand{\coord}{\mathsf{c}}            % coordinate
\newcommand{\prop}{\varphi}                 % proposition
\newcommand{\carrier}{\mathcal{A}}          % carrier type
\newcommand{\evid}{\mathcal{E}}             % evidence bundle
\newcommand{\oblig}{\mathcal{O}}            % obligations
\newcommand{\obstr}{\mathcal{B}}            % obstructions
\newcommand{\trust}{\mathcal{T}}            % trust annotation
\newcommand{\prov}{\Pi}                     % provenance

% Trust algebra
\newcommand{\Talg}{\mathfrak{T}}            % trust ordered algebra
\newcommand{\Eadm}{\mathcal{E}_{\mathrm{adm}}}
\newcommand{\tleq}{\preceq}                 % trust ordering
\newcommand{\tjoin}{\oplus}                 % conservative join
\newcommand{\tatten}{\ominus}               % attenuation
\newcommand{\tpromote}{\uparrow_\pi}        % promotion
\newcommand{\tdemote}{\downarrow_\chi}       % demotion

% Trust tiers
\newcommand{\Tcontra}{\ensuremath{\mathsf{CONTRADICTED}}}
\newcommand{\Tunver}{\ensuremath{\mathsf{UNVERIFIED}}}
\newcommand{\Tcopilot}{\ensuremath{\mathsf{COPILOT}}}
\newcommand{\Toracle}{\ensuremath{\mathsf{ORACLE}}}
\newcommand{\Truntime}{\ensuremath{\mathsf{RUNTIME}}}
\newcommand{\Tsolver}{\ensuremath{\mathsf{SOLVER}}}
\newcommand{\Tproof}{\ensuremath{\mathsf{PROOF}}}

% Site and geometry
\newcommand{\Site}{\mathcal{S}}             % semantic site
\newcommand{\Cov}{\mathrm{Cov}}            % covering family
\newcommand{\Ob}{\mathrm{Ob}}              % objects
\newcommand{\Mor}{\mathrm{Mor}}            % morphisms
\newcommand{\res}{\mathrm{res}}            % restriction
\newcommand{\Cech}{\check{C}}              % Čech complex
\newcommand{\Hcoh}[1]{H^{#1}}             % cohomology group

% Descent
\newcommand{\descent}{\mathrm{desc}}
\newcommand{\glue}{\mathrm{glue}}
\newcommand{\overlap}[2]{#1 \cap #2}

% Semantic moves
\newcommand{\VERIFY}{\textsc{Verify}}
\newcommand{\CONSTRUCT}{\textsc{Construct}}
\newcommand{\NORMALIZE}{\textsc{Normalize}}
\newcommand{\GLUE}{\textsc{Glue}}
\newcommand{\RESTRICT}{\textsc{Restrict}}
\newcommand{\TRANSPORT}{\textsc{Transport}}
\newcommand{\REPAIR}{\textsc{Repair}}
\newcommand{\PROMOTE}{\textsc{Promote}}
\newcommand{\CHALLENGE}{\textsc{Challenge}}
\newcommand{\ESCALATE}{\textsc{Escalate}}

% Fragments
\newcommand{\QFLIA}{\texttt{QF\_LIA}}
\newcommand{\QFLRA}{\texttt{QF\_LRA}}
\newcommand{\QFBV}{\texttt{QF\_BV}}
\newcommand{\QFUF}{\texttt{QF\_UF}}

% Misc
\newcommand{\Python}{\textsc{Python}}
\newcommand{\JuGeo}{\textsc{JuGeo}}
\newcommand{\LEAN}{\textsc{Lean}\,4}
\newcommand{\Fstar}{F$^{\star}$}
\newcommand{\Coq}{\textsc{Coq}}
\newcommand{\Dafny}{\textsc{Dafny}}

% Common shortcuts
\newcommand{\ie}{i.e.\@\xspace}
\newcommand{\eg}{e.g.\@\xspace}
\newcommand{\cf}{cf.\@\xspace}
\newcommand{\wrt}{w.r.t.\@\xspace}

% Evaluation shortcuts
\newcommand{\numcases}{300}
\newcommand{\numspec}{100}
\newcommand{\numeq}{100}
\newcommand{\numbug}{100}
\newcommand{\medtime}{6.5\,ms}
\newcommand{\totaltime}{1.949\,s}


% ── Packages ────────────────────────────────────────────────────────────
\usepackage{amsmath,amssymb,amsthm,mathtools}
\usepackage{tikz-cd}
\usepackage{listings}
\usepackage{xcolor}
\usepackage{xspace}
\usepackage{booktabs}
\usepackage{multirow}
\usepackage{enumitem}
\usepackage[expansion=false]{microtype}
\usepackage{hyperref}
\usepackage{cleveref}
% \usepackage{thmtools}  % disabled: not available in this TeX installation
\usepackage{float}
\usepackage{caption}
\usepackage{subcaption}
\usepackage{tabularx}
\usepackage{stmaryrd}       % \llbracket, \rrbracket
\usepackage{bbm}            % \mathbbm{1}

% ── Page geometry ───────────────────────────────────────────────────────
\usepackage[margin=1in]{geometry}

% ── Theorem environments ────────────────────────────────────────────────
\theoremstyle{plain}
\newtheorem{theorem}{Theorem}[section]
\newtheorem{lemma}[theorem]{Lemma}
\newtheorem{proposition}[theorem]{Proposition}
\newtheorem{corollary}[theorem]{Corollary}

\theoremstyle{definition}
\newtheorem{definition}[theorem]{Definition}
\newtheorem{example}[theorem]{Example}
\newtheorem{construction}[theorem]{Construction}

\theoremstyle{remark}
\newtheorem{remark}[theorem]{Remark}
\newtheorem{notation}[theorem]{Notation}

% ── Python listings style ──────────────────────────────────────────────
\definecolor{jugeoBlue}{HTML}{2563EB}
\definecolor{jugeoGreen}{HTML}{059669}
\definecolor{jugeoRed}{HTML}{DC2626}
\definecolor{jugeoPurple}{HTML}{7C3AED}
\definecolor{jugeoGray}{HTML}{6B7280}
\definecolor{jugeoBg}{HTML}{F8FAFC}

\lstdefinestyle{jugeo-python}{
  language=Python,
  basicstyle=\ttfamily\small,
  keywordstyle=\color{jugeoBlue}\bfseries,
  stringstyle=\color{jugeoGreen},
  commentstyle=\color{jugeoGray}\itshape,
  numberstyle=\tiny\color{jugeoGray},
  backgroundcolor=\color{jugeoBg},
  numbers=left,
  numbersep=6pt,
  frame=single,
  framerule=0.4pt,
  rulecolor=\color{jugeoGray!40},
  breaklines=true,
  breakatwhitespace=true,
  showstringspaces=false,
  tabsize=4,
  xleftmargin=1.5em,
  framexleftmargin=1.5em,
  morekeywords={dataclass,frozen,field,Enum,IntEnum,auto,Any,
                Mapping,Sequence,Optional,match,case},
  literate={→}{{$\to$}}1 {←}{{$\leftarrow$}}1
           {≤}{{$\leq$}}1 {≥}{{$\geq$}}1 {≠}{{$\neq$}}1
           {∈}{{$\in$}}1 {∀}{{$\forall$}}1 {∃}{{$\exists$}}1
           {⊕}{{$\oplus$}}1 {⊖}{{$\ominus$}}1,
  escapeinside={(*@}{@*)},
}
\lstset{style=jugeo-python}

\lstdefinestyle{jugeo-output}{
  basicstyle=\ttfamily\small,
  backgroundcolor=\color{jugeoBg},
  frame=single,
  framerule=0.4pt,
  rulecolor=\color{jugeoGray!40},
  breaklines=true,
  showstringspaces=false,
  xleftmargin=1.5em,
  framexleftmargin=0.5em,
  numbers=none,
}

% ── JuGeo-specific macros ──────────────────────────────────────────────

% Judgment 8-tuple
\newcommand{\J}{\mathcal{J}}
\newcommand{\Jud}[8]{(#1,\,#2,\,#3,\,#4,\,#5,\,#6,\,#7,\,#8)}
\newcommand{\JudShort}[1]{\J_{#1}}

% Components
\newcommand{\coord}{\mathsf{c}}            % coordinate
\newcommand{\prop}{\varphi}                 % proposition
\newcommand{\carrier}{\mathcal{A}}          % carrier type
\newcommand{\evid}{\mathcal{E}}             % evidence bundle
\newcommand{\oblig}{\mathcal{O}}            % obligations
\newcommand{\obstr}{\mathcal{B}}            % obstructions
\newcommand{\trust}{\mathcal{T}}            % trust annotation
\newcommand{\prov}{\Pi}                     % provenance

% Trust algebra
\newcommand{\Talg}{\mathfrak{T}}            % trust ordered algebra
\newcommand{\Eadm}{\mathcal{E}_{\mathrm{adm}}}
\newcommand{\tleq}{\preceq}                 % trust ordering
\newcommand{\tjoin}{\oplus}                 % conservative join
\newcommand{\tatten}{\ominus}               % attenuation
\newcommand{\tpromote}{\uparrow_\pi}        % promotion
\newcommand{\tdemote}{\downarrow_\chi}       % demotion

% Trust tiers
\newcommand{\Tcontra}{\ensuremath{\mathsf{CONTRADICTED}}}
\newcommand{\Tunver}{\ensuremath{\mathsf{UNVERIFIED}}}
\newcommand{\Tcopilot}{\ensuremath{\mathsf{COPILOT}}}
\newcommand{\Toracle}{\ensuremath{\mathsf{ORACLE}}}
\newcommand{\Truntime}{\ensuremath{\mathsf{RUNTIME}}}
\newcommand{\Tsolver}{\ensuremath{\mathsf{SOLVER}}}
\newcommand{\Tproof}{\ensuremath{\mathsf{PROOF}}}

% Site and geometry
\newcommand{\Site}{\mathcal{S}}             % semantic site
\newcommand{\Cov}{\mathrm{Cov}}            % covering family
\newcommand{\Ob}{\mathrm{Ob}}              % objects
\newcommand{\Mor}{\mathrm{Mor}}            % morphisms
\newcommand{\res}{\mathrm{res}}            % restriction
\newcommand{\Cech}{\check{C}}              % Čech complex
\newcommand{\Hcoh}[1]{H^{#1}}             % cohomology group

% Descent
\newcommand{\descent}{\mathrm{desc}}
\newcommand{\glue}{\mathrm{glue}}
\newcommand{\overlap}[2]{#1 \cap #2}

% Semantic moves
\newcommand{\VERIFY}{\textsc{Verify}}
\newcommand{\CONSTRUCT}{\textsc{Construct}}
\newcommand{\NORMALIZE}{\textsc{Normalize}}
\newcommand{\GLUE}{\textsc{Glue}}
\newcommand{\RESTRICT}{\textsc{Restrict}}
\newcommand{\TRANSPORT}{\textsc{Transport}}
\newcommand{\REPAIR}{\textsc{Repair}}
\newcommand{\PROMOTE}{\textsc{Promote}}
\newcommand{\CHALLENGE}{\textsc{Challenge}}
\newcommand{\ESCALATE}{\textsc{Escalate}}

% Fragments
\newcommand{\QFLIA}{\texttt{QF\_LIA}}
\newcommand{\QFLRA}{\texttt{QF\_LRA}}
\newcommand{\QFBV}{\texttt{QF\_BV}}
\newcommand{\QFUF}{\texttt{QF\_UF}}

% Misc
\newcommand{\Python}{\textsc{Python}}
\newcommand{\JuGeo}{\textsc{JuGeo}}
\newcommand{\LEAN}{\textsc{Lean}\,4}
\newcommand{\Fstar}{F$^{\star}$}
\newcommand{\Coq}{\textsc{Coq}}
\newcommand{\Dafny}{\textsc{Dafny}}

% Common shortcuts
\newcommand{\ie}{i.e.\@\xspace}
\newcommand{\eg}{e.g.\@\xspace}
\newcommand{\cf}{cf.\@\xspace}
\newcommand{\wrt}{w.r.t.\@\xspace}

% Evaluation shortcuts
\newcommand{\numcases}{300}
\newcommand{\numspec}{100}
\newcommand{\numeq}{100}
\newcommand{\numbug}{100}
\newcommand{\medtime}{6.5\,ms}
\newcommand{\totaltime}{1.949\,s}


% ── Packages ────────────────────────────────────────────────────────────
\usepackage{amsmath,amssymb,amsthm,mathtools}
\usepackage{tikz-cd}
\usepackage{listings}
\usepackage{xcolor}
\usepackage{xspace}
\usepackage{booktabs}
\usepackage{multirow}
\usepackage{enumitem}
\usepackage[expansion=false]{microtype}
\usepackage{hyperref}
\usepackage{cleveref}
% \usepackage{thmtools}  % disabled: not available in this TeX installation
\usepackage{float}
\usepackage{caption}
\usepackage{subcaption}
\usepackage{tabularx}
\usepackage{stmaryrd}       % \llbracket, \rrbracket
\usepackage{bbm}            % \mathbbm{1}

% ── Page geometry ───────────────────────────────────────────────────────
\usepackage[margin=1in]{geometry}

% ── Theorem environments ────────────────────────────────────────────────
\theoremstyle{plain}
\newtheorem{theorem}{Theorem}[section]
\newtheorem{lemma}[theorem]{Lemma}
\newtheorem{proposition}[theorem]{Proposition}
\newtheorem{corollary}[theorem]{Corollary}

\theoremstyle{definition}
\newtheorem{definition}[theorem]{Definition}
\newtheorem{example}[theorem]{Example}
\newtheorem{construction}[theorem]{Construction}

\theoremstyle{remark}
\newtheorem{remark}[theorem]{Remark}
\newtheorem{notation}[theorem]{Notation}

% ── Python listings style ──────────────────────────────────────────────
\definecolor{jugeoBlue}{HTML}{2563EB}
\definecolor{jugeoGreen}{HTML}{059669}
\definecolor{jugeoRed}{HTML}{DC2626}
\definecolor{jugeoPurple}{HTML}{7C3AED}
\definecolor{jugeoGray}{HTML}{6B7280}
\definecolor{jugeoBg}{HTML}{F8FAFC}

\lstdefinestyle{jugeo-python}{
  language=Python,
  basicstyle=\ttfamily\small,
  keywordstyle=\color{jugeoBlue}\bfseries,
  stringstyle=\color{jugeoGreen},
  commentstyle=\color{jugeoGray}\itshape,
  numberstyle=\tiny\color{jugeoGray},
  backgroundcolor=\color{jugeoBg},
  numbers=left,
  numbersep=6pt,
  frame=single,
  framerule=0.4pt,
  rulecolor=\color{jugeoGray!40},
  breaklines=true,
  breakatwhitespace=true,
  showstringspaces=false,
  tabsize=4,
  xleftmargin=1.5em,
  framexleftmargin=1.5em,
  morekeywords={dataclass,frozen,field,Enum,IntEnum,auto,Any,
                Mapping,Sequence,Optional,match,case},
  literate={→}{{$\to$}}1 {←}{{$\leftarrow$}}1
           {≤}{{$\leq$}}1 {≥}{{$\geq$}}1 {≠}{{$\neq$}}1
           {∈}{{$\in$}}1 {∀}{{$\forall$}}1 {∃}{{$\exists$}}1
           {⊕}{{$\oplus$}}1 {⊖}{{$\ominus$}}1,
  escapeinside={(*@}{@*)},
}
\lstset{style=jugeo-python}

\lstdefinestyle{jugeo-output}{
  basicstyle=\ttfamily\small,
  backgroundcolor=\color{jugeoBg},
  frame=single,
  framerule=0.4pt,
  rulecolor=\color{jugeoGray!40},
  breaklines=true,
  showstringspaces=false,
  xleftmargin=1.5em,
  framexleftmargin=0.5em,
  numbers=none,
}

% ── JuGeo-specific macros ──────────────────────────────────────────────

% Judgment 8-tuple
\newcommand{\J}{\mathcal{J}}
\newcommand{\Jud}[8]{(#1,\,#2,\,#3,\,#4,\,#5,\,#6,\,#7,\,#8)}
\newcommand{\JudShort}[1]{\J_{#1}}

% Components
\newcommand{\coord}{\mathsf{c}}            % coordinate
\newcommand{\prop}{\varphi}                 % proposition
\newcommand{\carrier}{\mathcal{A}}          % carrier type
\newcommand{\evid}{\mathcal{E}}             % evidence bundle
\newcommand{\oblig}{\mathcal{O}}            % obligations
\newcommand{\obstr}{\mathcal{B}}            % obstructions
\newcommand{\trust}{\mathcal{T}}            % trust annotation
\newcommand{\prov}{\Pi}                     % provenance

% Trust algebra
\newcommand{\Talg}{\mathfrak{T}}            % trust ordered algebra
\newcommand{\Eadm}{\mathcal{E}_{\mathrm{adm}}}
\newcommand{\tleq}{\preceq}                 % trust ordering
\newcommand{\tjoin}{\oplus}                 % conservative join
\newcommand{\tatten}{\ominus}               % attenuation
\newcommand{\tpromote}{\uparrow_\pi}        % promotion
\newcommand{\tdemote}{\downarrow_\chi}       % demotion

% Trust tiers
\newcommand{\Tcontra}{\ensuremath{\mathsf{CONTRADICTED}}}
\newcommand{\Tunver}{\ensuremath{\mathsf{UNVERIFIED}}}
\newcommand{\Tcopilot}{\ensuremath{\mathsf{COPILOT}}}
\newcommand{\Toracle}{\ensuremath{\mathsf{ORACLE}}}
\newcommand{\Truntime}{\ensuremath{\mathsf{RUNTIME}}}
\newcommand{\Tsolver}{\ensuremath{\mathsf{SOLVER}}}
\newcommand{\Tproof}{\ensuremath{\mathsf{PROOF}}}

% Site and geometry
\newcommand{\Site}{\mathcal{S}}             % semantic site
\newcommand{\Cov}{\mathrm{Cov}}            % covering family
\newcommand{\Ob}{\mathrm{Ob}}              % objects
\newcommand{\Mor}{\mathrm{Mor}}            % morphisms
\newcommand{\res}{\mathrm{res}}            % restriction
\newcommand{\Cech}{\check{C}}              % Čech complex
\newcommand{\Hcoh}[1]{H^{#1}}             % cohomology group

% Descent
\newcommand{\descent}{\mathrm{desc}}
\newcommand{\glue}{\mathrm{glue}}
\newcommand{\overlap}[2]{#1 \cap #2}

% Semantic moves
\newcommand{\VERIFY}{\textsc{Verify}}
\newcommand{\CONSTRUCT}{\textsc{Construct}}
\newcommand{\NORMALIZE}{\textsc{Normalize}}
\newcommand{\GLUE}{\textsc{Glue}}
\newcommand{\RESTRICT}{\textsc{Restrict}}
\newcommand{\TRANSPORT}{\textsc{Transport}}
\newcommand{\REPAIR}{\textsc{Repair}}
\newcommand{\PROMOTE}{\textsc{Promote}}
\newcommand{\CHALLENGE}{\textsc{Challenge}}
\newcommand{\ESCALATE}{\textsc{Escalate}}

% Fragments
\newcommand{\QFLIA}{\texttt{QF\_LIA}}
\newcommand{\QFLRA}{\texttt{QF\_LRA}}
\newcommand{\QFBV}{\texttt{QF\_BV}}
\newcommand{\QFUF}{\texttt{QF\_UF}}

% Misc
\newcommand{\Python}{\textsc{Python}}
\newcommand{\JuGeo}{\textsc{JuGeo}}
\newcommand{\LEAN}{\textsc{Lean}\,4}
\newcommand{\Fstar}{F$^{\star}$}
\newcommand{\Coq}{\textsc{Coq}}
\newcommand{\Dafny}{\textsc{Dafny}}

% Common shortcuts
\newcommand{\ie}{i.e.\@\xspace}
\newcommand{\eg}{e.g.\@\xspace}
\newcommand{\cf}{cf.\@\xspace}
\newcommand{\wrt}{w.r.t.\@\xspace}

% Evaluation shortcuts
\newcommand{\numcases}{300}
\newcommand{\numspec}{100}
\newcommand{\numeq}{100}
\newcommand{\numbug}{100}
\newcommand{\medtime}{6.5\,ms}
\newcommand{\totaltime}{1.949\,s}

% experiment-data.tex — AUTO-GENERATED from real jugeo CLI runs
% DO NOT EDIT — regenerate with: python3 experiments/run_universal_metrics.py
% Generated from 30 programs across 6 domains

\newcommand{\expTotalPrograms}{30}
\newcommand{\expVerified}{30}
\newcommand{\expAccuracy}{100.0\%}
\newcommand{\expCoordsMin}{2}
\newcommand{\expCoordsMax}{10}
\newcommand{\expCoordsMean}{5.2}
\newcommand{\expCoordsMedian}{5.5}
\newcommand{\expPropsMin}{4}
\newcommand{\expPropsMax}{20}
\newcommand{\expPropsMean}{10.4}
\newcommand{\expPropsSum}{311}
\newcommand{\expPropsOkSum}{310}
\newcommand{\expTimeMin}{3.506\,s}
\newcommand{\expTimeMax}{6.714\,s}
\newcommand{\expTimeMean}{4.64\,s}
\newcommand{\expTimeMedian}{4.508\,s}
\newcommand{\expTimeTotal}{139.204\,s}
\newcommand{\expObstructionTotal}{0}
\newcommand{\expHOneZero}{30}
\newcommand{\expDomainCount}{6}
\newcommand{\expTrustCOPILOTSUGGESTED}{30}
\newcommand{\expDomainDsCount}{5}
\newcommand{\expDomainDsVerified}{5}
\newcommand{\expDomainDsAvgCoords}{7.6}
\newcommand{\expDomainDsAvgProps}{15.2}
\newcommand{\expDomainMathCount}{5}
\newcommand{\expDomainMathVerified}{5}
\newcommand{\expDomainMathAvgCoords}{4.8}
\newcommand{\expDomainMathAvgProps}{9.4}
\newcommand{\expDomainSmCount}{5}
\newcommand{\expDomainSmVerified}{5}
\newcommand{\expDomainSmAvgCoords}{7.6}
\newcommand{\expDomainSmAvgProps}{15.2}
\newcommand{\expDomainSortCount}{5}
\newcommand{\expDomainSortVerified}{5}
\newcommand{\expDomainSortAvgCoords}{2.4}
\newcommand{\expDomainSortAvgProps}{4.8}
\newcommand{\expDomainStrCount}{5}
\newcommand{\expDomainStrVerified}{5}
\newcommand{\expDomainStrAvgCoords}{3.2}
\newcommand{\expDomainStrAvgProps}{6.4}
\newcommand{\expDomainWebCount}{5}
\newcommand{\expDomainWebVerified}{5}
\newcommand{\expDomainWebAvgCoords}{5.6}
\newcommand{\expDomainWebAvgProps}{11.2}

% subsystem-data.tex — AUTO-GENERATED from real jugeo subsystem experiments
% DO NOT EDIT — regenerate with: python3 experiments/run_subsystem_experiments.py

\newcommand{\subCliTotal}{10}
\newcommand{\subCliVerified}{9}
\newcommand{\subCliAccuracy}{90.0\%}
\newcommand{\subCliMeanTime}{4.132\,s}
\newcommand{\subCliMinTime}{3.472\,s}
\newcommand{\subCliMaxTime}{5.372\,s}
\newcommand{\subCliTotalCoords}{54}
\newcommand{\subCliTotalProps}{107}
\newcommand{\subCliTotalPropsOk}{106}
\newcommand{\subSiteCalls}{90}
\newcommand{\subSiteSuccessful}{69}
\newcommand{\subSiteSuccessRate}{76.7\%}
\newcommand{\subSiteMeanTime}{0.0001\,s}
\newcommand{\subSiteTotalTime}{0.005\,s}
\newcommand{\subSpecificationsatisfactionSmallTime}{0.0003\,s}
\newcommand{\subSpecificationsatisfactionMediumTime}{0.0001\,s}
\newcommand{\subSpecificationsatisfactionLargeTime}{0.0001\,s}
\newcommand{\subInhabitantfleetSmallTime}{0.0\,s}
\newcommand{\subInhabitantfleetMediumTime}{0.0\,s}
\newcommand{\subInhabitantfleetLargeTime}{0.0\,s}
\newcommand{\subTheoremecologySmallTime}{0.0\,s}
\newcommand{\subTheoremecologyMediumTime}{0.0\,s}
\newcommand{\subTheoremecologyLargeTime}{0.0\,s}
\newcommand{\subStatespaceexplorationSmallTime}{0.0\,s}
\newcommand{\subStatespaceexplorationMediumTime}{0.0\,s}
\newcommand{\subStatespaceexplorationLargeTime}{0.0\,s}
\newcommand{\subRepairsemanticsSmallTime}{0.0\,s}
\newcommand{\subRepairsemanticsMediumTime}{0.0\,s}
\newcommand{\subRepairsemanticsLargeTime}{0.0\,s}
\newcommand{\subReplaygluingSmallTime}{0.0\,s}
\newcommand{\subReplaygluingMediumTime}{0.0\,s}
\newcommand{\subReplaygluingLargeTime}{0.0\,s}
\newcommand{\subSemanticfuturesSmallTime}{0.0\,s}
\newcommand{\subSemanticfuturesMediumTime}{0.0\,s}
\newcommand{\subSemanticfuturesLargeTime}{0.0\,s}
\newcommand{\subHypercovertreatySmallTime}{0.0\,s}
\newcommand{\subHypercovertreatyMediumTime}{0.0\,s}
\newcommand{\subHypercovertreatyLargeTime}{0.0\,s}
\newcommand{\subEncodeforsolverSmallTime}{0.0026\,s}
\newcommand{\subEncodeforsolverMediumTime}{0.0002\,s}
\newcommand{\subEncodeforsolverLargeTime}{0.0001\,s}
\newcommand{\subInterfaceroutingSmallTime}{0.0\,s}
\newcommand{\subInterfaceroutingMediumTime}{0.0\,s}
\newcommand{\subInterfaceroutingLargeTime}{0.0\,s}
\newcommand{\subOrchestrateverificationSmallTime}{0.0002\,s}
\newcommand{\subOrchestrateverificationMediumTime}{0.0\,s}
\newcommand{\subOrchestrateverificationLargeTime}{0.0\,s}
\newcommand{\subRunfulldescentSmallTime}{0.0\,s}
\newcommand{\subRunfulldescentMediumTime}{0.0\,s}
\newcommand{\subRunfulldescentLargeTime}{0.0\,s}
\newcommand{\subKernellifecycleSmallTime}{0.0\,s}
\newcommand{\subKernellifecycleMediumTime}{0.0\,s}
\newcommand{\subKernellifecycleLargeTime}{0.0\,s}
\newcommand{\subDiscoverypipelineSmallTime}{0.0\,s}
\newcommand{\subDiscoverypipelineMediumTime}{0.0\,s}
\newcommand{\subDiscoverypipelineLargeTime}{0.0\,s}
\newcommand{\subAnalogytransportSmallTime}{0.0\,s}
\newcommand{\subAnalogytransportMediumTime}{0.0\,s}
\newcommand{\subAnalogytransportLargeTime}{0.0\,s}
\newcommand{\subFormalcoresiteSmallTime}{0.0001\,s}
\newcommand{\subFormalcoresiteMediumTime}{0.0\,s}
\newcommand{\subFormalcoresiteLargeTime}{0.0\,s}
\newcommand{\subProblematlasSmallTime}{0.0\,s}
\newcommand{\subProblematlasMediumTime}{0.0\,s}
\newcommand{\subProblematlasLargeTime}{0.0\,s}
\newcommand{\subPublicalignmentSmallTime}{0.0\,s}
\newcommand{\subPublicalignmentMediumTime}{0.0\,s}
\newcommand{\subPublicalignmentLargeTime}{0.0\,s}
\newcommand{\subRegimebootstrappingSmallTime}{0.0\,s}
\newcommand{\subRegimebootstrappingMediumTime}{0.0\,s}
\newcommand{\subRegimebootstrappingLargeTime}{0.0\,s}
\newcommand{\subRelationalrefinementSmallTime}{0.0\,s}
\newcommand{\subRelationalrefinementMediumTime}{0.0\,s}
\newcommand{\subRelationalrefinementLargeTime}{0.0\,s}
\newcommand{\subSemanticclosureSmallTime}{0.0\,s}
\newcommand{\subSemanticclosureMediumTime}{0.0\,s}
\newcommand{\subSemanticclosureLargeTime}{0.0\,s}
\newcommand{\subTheoremeconomicsSmallTime}{0.0001\,s}
\newcommand{\subTheoremeconomicsMediumTime}{0.0\,s}
\newcommand{\subTheoremeconomicsLargeTime}{0.0\,s}
\newcommand{\subBenchmarksuiteSmallTime}{0.0\,s}
\newcommand{\subBenchmarksuiteMediumTime}{0.0\,s}
\newcommand{\subBenchmarksuiteLargeTime}{0.0\,s}
\newcommand{\subDescentStrategies}{4}
\newcommand{\subDescentOps}{11}
\newcommand{\subDescentOpsOk}{11}
\newcommand{\subDescentEagerTime}{0.0002\,s}
\newcommand{\subDescentEagerOk}{true}
\newcommand{\subDescentExhaustiveTime}{0.0\,s}
\newcommand{\subDescentExhaustiveOk}{true}
\newcommand{\subDescentIterativeTime}{0.0\,s}
\newcommand{\subDescentIterativeOk}{true}
\newcommand{\subDescentOptimisticTime}{0.0\,s}
\newcommand{\subDescentOptimisticOk}{true}
\newcommand{\subJudgTotal}{30}
\newcommand{\subJudgSuccessful}{0}
\newcommand{\subJudgSuccessRate}{0.0\%}
\newcommand{\subJudgMeanTimeUs}{7.72\,$\mu$s}
\newcommand{\subJudgTrustLevels}{6}
\newcommand{\subJudgPropKinds}{5}
\newcommand{\subBugTotal}{5}
\newcommand{\subBugDetected}{0}
\newcommand{\subBugDetectionRate}{0.0\%}
\newcommand{\subBugFalsePositiveRate}{0.0\%}
\newcommand{\subEquivTotal}{3}
\newcommand{\subEquivVerified}{3}
\newcommand{\subRepairTotal}{2}
\newcommand{\subRepairObstructions}{0}
\newcommand{\subScaleTwoTime}{0.0001\,s}
\newcommand{\subScaleFourTime}{0.0\,s}
\newcommand{\subScaleEightTime}{0.0\,s}
\newcommand{\subScaleTwelveTime}{0.0\,s}
\newcommand{\subScaleSixteenTime}{0.0\,s}
\newcommand{\subScaleTwentyTime}{0.0\,s}
\newcommand{\subTotalExperimentTime}{95.6\,s}

% data-paper45.tex -- AUTO-GENERATED by exp45_callable_surfaces.py
% DO NOT EDIT -- regenerate with: python3 experiments/exp45_callable_surfaces.py

\newcommand{\ppFortyFiveTotalPrograms}{10}
\newcommand{\ppFortyFiveNumHof}{4}
\newcommand{\ppFortyFiveNumCb}{3}
\newcommand{\ppFortyFiveNumDec}{3}

% --- Recall (fraction verified) per category ---
\newcommand{\ppFortyFiveHofRecall}{100.0\%}
\newcommand{\ppFortyFiveCbRecall}{100.0\%}
\newcommand{\ppFortyFiveDecRecall}{100.0\%}

% --- Coverage (mean evaluate score) per category (replaces --- in table) ---
\newcommand{\ppFortyFiveHofCoverage}{0.0\%}
\newcommand{\ppFortyFiveCbCoverage}{0.0\%}
\newcommand{\ppFortyFiveDecCoverage}{0.0\%}

% --- Timing ---
\newcommand{\ppFortyFiveMedTime}{5.90\,s}

% --- Surfaces ---
\newcommand{\ppFortyFiveMeanSurfaces}{3.8}


% ── Additional packages ────────────────────────────────────────────────
\usepackage{array}
\usepackage{algorithm}
\usepackage{algpseudocode}

% ── Paper-specific macros ──────────────────────────────────────────────
\newcommand{\CS}{\mathcal{CS}}                 % callable surface (set)
\newcommand{\surf}{\sigma}                     % individual surface
\newcommand{\Surf}{\Sigma}                     % surface space
\newcommand{\CEnv}{\Gamma}                     % contract environment
\newcommand{\Clos}{\mathsf{Clos}}              % closure
\newcommand{\FV}{\mathrm{FV}}                  % free variables
\newcommand{\BV}{\mathrm{BV}}                  % bound variables
\newcommand{\Env}{\mathcal{E}}                 % value environment
\newcommand{\Decor}{\mathsf{Dec}}              % decorator
\newcommand{\CB}{\mathsf{CB}}                  % callback
\newcommand{\Chain}{\mathsf{Chain}}            % callback chain
\newcommand{\HOF}{\mathsf{HOF}}                % higher-order function
\newcommand{\Contract}{\mathsf{Con}}           % contract
\newcommand{\Pre}{\mathsf{pre}}                % precondition
\newcommand{\Post}{\mathsf{post}}              % postcondition
\newcommand{\sound}{\mathsf{sound}}            % soundness predicate
\newcommand{\complete}{\mathsf{cmplt}}         % completeness predicate
\newcommand{\reach}{\mathrm{reach}}            % reachable surfaces
\newcommand{\compose}{\mathbin{\odot}}         % surface composition
\newcommand{\wrap}{\mathsf{wrap}}              % decorator wrapping
\newcommand{\capture}{\mathsf{cap}}            % captured-variable map
\newcommand{\SiteMor}{\mathsf{SMor}}           % site morphism
\newcommand{\ContEnv}{\mathsf{CE}}             % contract environment functor
\newcommand{\MapF}{\texttt{map}}               % Python map
\newcommand{\FilterF}{\texttt{filter}}         % Python filter
\newcommand{\ReduceF}{\texttt{reduce}}         % Python reduce
\newcommand{\HOSound}{\textbf{HO-Sound}}       % higher-order soundness label
\newcommand{\CallSurf}{\textsc{CallSurf}}      % system name
\newcommand{\numhof}{87}                       % benchmark HOF patterns
\newcommand{\numcb}{54}                        % benchmark callback patterns
\newcommand{\hofrecall}{94.3}                  % HOF recall %
\newcommand{\cbrecall}{91.8}                   % callback recall %

\title{\textbf{Callable Surface Analysis:\\
       Verifying Higher-Order \Python{} Functions\\
       via Morphisms in the Semantic Site}}

\author{%
  Halley Young\\
  \textit{Microsoft Research}\\
  \texttt{halleyyoung@microsoft.com}
}

\date{\today}

\begin{document}
\maketitle

% =====================================================================
\begin{abstract}
Higher-order functions---closures, decorators, callbacks, and the
classic \MapF/\FilterF/\ReduceF{} trio---are ubiquitous in idiomatic
\Python{} yet impose formidable challenges for compositional
verification.  The difficulty is structural: control crosses function
boundaries through first-class callable values that are only known at
run time, erasing the static call-graph edges on which most existing
analyses rely.

We introduce \emph{callable surface analysis}, a framework that
identifies every boundary where control flow crosses a
function boundary---a \emph{callable surface}---and maps each
such surface to a morphism in the \JuGeo{} semantic site
$(\Site, J, \Cov)$.  A callable surface $\surf$ carries a
\emph{contract} $\Contract(\surf) = (\Pre, \Post, \tau)$ whose
trust level $\tau \in \Talg$ is propagated from the site morphism.
Three specialised sub-analyses handle the principal surface
families: (i)~\emph{closure analysis} tracks captured variables
through $\lambda$-lifting to produce a closed contract for each
closure; (ii)~\emph{decorator verification} models a decorator as
morphism composition $f \compose g$ in the site and derives the
composite contract algebraically; (iii)~\emph{callback-chain
verification} treats a chain of callbacks as a sequence of
restricted morphisms and checks contract compatibility at each
link.

Our central theoretical result---the \emph{Higher-Order Soundness
Theorem}---states that if every callable surface in a
\Python{} program is equipped with a verified contract, then
higher-order composition cannot produce a contract escape: no
unchecked callable value can be invoked without passing through a
verified morphism in the site.  We prove the theorem formally in
\LEAN{} with a complete, sorry-free mechanisation.

We implement the framework in
\texttt{src/jugeo/python\_runtime/callable\_surfaces/} and evaluate
it on a benchmark of $\numcases{}$ \Python{} programs drawn from
real-world open-source repositories.  On $\numhof{}$ higher-order
function patterns, \CallSurf{} achieves $\hofrecall{}\,\%$ recall;
on $\numcb{}$ callback patterns it achieves $\cbrecall{}\,\%$
recall.  Median per-site construction time is $\medtime{}$.
\end{abstract}

\newpage

% =====================================================================
\section{Introduction}\label{sec:intro}
% =====================================================================

\Python{} is one of the most widely deployed programming languages in
the world.  Its success stems, in part, from the first-class status of
callable objects: functions are values, functions can be passed to
other functions, functions can be returned from functions, and
functions can be constructed at run time through closures and
decorators.  This expressiveness is a virtue for rapid development,
but it creates a fundamental challenge for verification: the control
graph is \emph{not} statically determined.

\paragraph{The callable-boundary problem.}
Consider the canonical pattern:
\begin{lstlisting}[style=jugeo-python]
results = list(map(validate, user_inputs))
\end{lstlisting}
The contract of \texttt{results} depends on the contract of
\texttt{validate}, which is passed as a \emph{first-class value}.
A conventional Hoare-logic analyser processes \texttt{map} without
knowing what \texttt{validate} does; a whole-program analyser must
inline every possible callee, which is expensive and sometimes
undecidable.  Neither approach scales to the decorator-heavy,
callback-driven code prevalent in modern Python frameworks such as
Flask, FastAPI, and asyncio.

\paragraph{Our approach: callable surfaces.}
We model the problem geometrically.  Every boundary where control
crosses a function boundary corresponds to a \emph{callable surface}
$\surf \in \CS$ in the abstract syntax tree (AST).  A callable
surface is not a call site in the classical sense; it is a
\emph{potential crossing point}---a place where a callable value
might be invoked.  Mapping callable surfaces to morphisms in the
\JuGeo{} semantic site $(\Site, J, \Cov)$ lets us:
\begin{enumerate}
  \item attach a \emph{contract} $\Contract(\surf)$ to each surface
        via the judgment algebra;
  \item propagate trust levels along morphism composition;
  \item verify higher-order patterns by checking contract
        compatibility at each surface.
\end{enumerate}

\paragraph{Main contributions.}
This paper contributes:
\begin{enumerate}
  \item The \emph{callable surface model} (\S\ref{sec:model}):
        a taxonomy of surface kinds (direct, closure, decorator,
        callback, HOF-argument) and a canonical AST extraction
        algorithm.
  \item \emph{Closure analysis} (\S\ref{sec:closure}):
        a $\lambda$-lifting pass that makes captured variables
        explicit, enabling closed contracts for closures.
  \item \emph{Decorator verification} (\S\ref{sec:decorator}):
        a treatment of decorators as morphism composition in the site,
        with an algebraic contract derivation rule.
  \item \emph{Callback-chain verification} (\S\ref{sec:callbacks}):
        a sequential morphism model for event-driven and promise-style
        code.
  \item The \emph{site construction} procedure for callable surfaces
        (\S\ref{sec:site}): a functor from the callable surface
        category to the semantic site.
  \item The \emph{Higher-Order Soundness Theorem}
        (\S\ref{sec:theorem}): a formally verified guarantee that no
        callable escape is possible under complete surface coverage,
        with a sorry-free \LEAN{} proof.
  \item \emph{Experimental evaluation} (\S\ref{sec:experiments}) on
        $\numcases{}$ programs.
\end{enumerate}

\paragraph{Organisation.}
Sections \ref{sec:model}--\ref{sec:site} develop the theory;
\S\ref{sec:theorem} states and proves the main theorem;
\S\ref{sec:experiments} reports experiments; \S\ref{sec:conclusion}
concludes.

% =====================================================================
\section{Callable Surface Model}\label{sec:model}
% =====================================================================

\subsection{Surface Kinds}

We distinguish five \emph{surface kinds} based on how a callable
value is invoked.

\begin{definition}[Surface Kind]\label{def:surfkind}
A \emph{surface kind} $k \in \{\mathsf{direct},\, \mathsf{closure},\,
\mathsf{decorator},\, \mathsf{callback},\, \mathsf{hof\_arg}\}$
classifies a callable surface by the syntactic form of its invocation:
\begin{itemize}
  \item $\mathsf{direct}$: a call to a named function,
        \texttt{f(x)}.
  \item $\mathsf{closure}$: a call through a variable bound to a
        closure, \texttt{g(x)} where \texttt{g} captures variables
        from an enclosing scope.
  \item $\mathsf{decorator}$: an application of a decorator
        \texttt{@d} to a function definition.
  \item $\mathsf{callback}$: a call to a function passed as an
        event handler or promise continuation.
  \item $\mathsf{hof\_arg}$: a call to a function passed as an
        argument to a higher-order function such as \MapF{},
        \FilterF{}, or \ReduceF{}.
\end{itemize}
\end{definition}

\begin{definition}[Callable Surface]\label{def:surface}
A \emph{callable surface} is a tuple
$\surf = (\mathit{id},\, k,\, n,\, \mu)$ where
$\mathit{id}$ is a unique identifier,
$k$ is the surface kind (Definition~\ref{def:surfkind}),
$n$ is the arity of the callable,
and $\mu \in \mathbb{N}$ is the capture count (number of free
variables captured by closures; $0$ for non-closure kinds).
\end{definition}

\subsection{AST Extraction Algorithm}

The extraction algorithm walks the AST of a \Python{} module and
emits a callable surface at each syntactic location where a callable
value may be invoked.

\begin{algorithm}[H]
\caption{$\textsc{ExtractSurfaces}(\mathit{ast})$}
\begin{algorithmic}[1]
\State $\CS \gets \emptyset$
\ForAll{node $n$ in $\mathrm{pre\text{-}order}(\mathit{ast})$}
  \If{$n$ is \texttt{Call}$(f, \mathit{args})$}
    \State $k \gets \mathsf{classifyCallee}(f)$
    \State $\surf \gets (\mathrm{fresh}(),\, k,\, |\mathit{args}|,\,
            \mathrm{captureCount}(f))$
    \State $\CS \gets \CS \cup \{\surf\}$
  \EndIf
  \If{$n$ is \texttt{FunctionDef} with decorators $\mathit{decs}$}
    \ForAll{$d \in \mathit{decs}$}
      \State $\surf \gets (\mathrm{fresh}(),\, \mathsf{decorator},\,
              1,\, 0)$
      \State $\CS \gets \CS \cup \{\surf\}$
    \EndFor
  \EndIf
\EndFor
\Return $\CS$
\end{algorithmic}
\end{algorithm}

The helper $\mathsf{classifyCallee}$ inspects the AST node of the
callee: a \texttt{Name} node whose binding is a free variable in the
enclosing scope maps to $\mathsf{closure}$; a \texttt{Name} passed as
an argument to a known HOF maps to $\mathsf{hof\_arg}$; an event
handler registration maps to $\mathsf{callback}$; otherwise the
surface is $\mathsf{direct}$.

\begin{proposition}[Extraction Completeness]\label{prop:extract}
For any \Python{} program $P$ in the Core Python fragment (no
\texttt{eval}, no metaclasses), $\textsc{ExtractSurfaces}(\mathit{ast}(P))$
returns a set $\CS$ such that every dynamic call site of $P$ is
represented by at least one surface $\surf \in \CS$.
\end{proposition}

\begin{proof}[Proof sketch]
By structural induction on the AST.  The pre-order traversal visits
every node.  Every dynamic call is syntactically represented by a
\texttt{Call} node; every decorator application by a
\texttt{FunctionDef} node with non-empty decorator list.  The
completeness of the cover follows from the completeness of the AST
representation of Core Python, which excludes only
\texttt{eval}-generated code.
\end{proof}

\subsection{Contract Assignment}

Each surface receives a contract from the judgment algebra.

\begin{definition}[Surface Contract]\label{def:contract}
A \emph{surface contract} for $\surf$ is a triple
$\Contract(\surf) = (\Pre, \Post, \tau)$ where
$\Pre, \Post$ are propositions over the carrier type of the
associated judgment, and $\tau \in \Talg$ is the trust level.
The \emph{contract environment} $\CEnv : \CS \to \Contract$
assigns a contract to every surface.
\end{definition}

A surface is \emph{contractually sound} if its contract is verified:
$\sound(\surf, \CEnv) \iff \CEnv(\surf).\tau \ge \Tsolver$.

% =====================================================================
\section{Closure Analysis}\label{sec:closure}
% =====================================================================

Closures are \Python's primary mechanism for constructing first-class
functions with enclosed state.  A closure $c$ captures a set of
free variables $\FV(c) \subseteq \mathrm{Vars}$ from its defining
scope.  To verify a closure we must account for the captured
bindings: an erroneous value in the captured environment can
invalidate a contract even if the closure body itself is correct.

\subsection{Lambda-Lifting Transform}

We apply a standard $\lambda$-lifting (closure conversion) transform
that makes captured variables explicit as additional parameters.

\begin{definition}[$\lambda$-Lift]\label{def:lambdalift}
Let $c$ be a closure with body $b$ and captured variables
$\vec{v} = (v_1,\dots,v_m) = \FV(c)$.
The \emph{$\lambda$-lift} of $c$ is a top-level function
$\hat{c}(\vec{x}, \vec{v})$ where $\vec{x}$ are the original
parameters and $\vec{v}$ are the captured variables, now passed
explicitly.
\end{definition}

\begin{construction}[Closure Contract Derivation]\label{const:closcon}
Given a closure $c$ with environment $\Env$ assigning trust levels to
captured variables, the closed contract for $c$ is:
\[
  \Contract(c) = \Bigl(\Pre(c) \wedge \bigwedge_{v \in \FV(c)} \tau_\Env(v) \ge \tau_{\min},\;
                       \Post(c),\;
                       \bigwedge_{v \in \FV(c)} \tau_\Env(v) \Bigr)
\]
where $\tau_{\min}$ is the minimum trust level required by the closure
body and $\tau_\Env(v)$ is the trust level of the current binding of $v$.
\end{construction}

\subsection{Captured-Variable Tracking}

The analyzer maintains a \emph{capture map} $\capture :
\mathrm{Closures} \to \mathcal{P}(\mathrm{Coordinates})$ that maps
each closure to the set of coordinates whose judgments must be
consulted when verifying the closure's contract.

\begin{lemma}[Closure Contract Soundness]\label{lem:closuresound}
Let $c$ be a closure with captured variables $\FV(c)$.  If
\begin{enumerate}
  \item the $\lambda$-lifted body $\hat{c}$ satisfies its contract
        under the extended parameter list, and
  \item for every $v \in \FV(c)$, the judgment at coordinate
        $\capture(c)(v)$ has trust level $\ge \tau_{\min}$,
\end{enumerate}
then $c$ satisfies $\Contract(c)$ in the original environment.
\end{lemma}

\begin{proof}
By the soundness of $\lambda$-lifting (see~\cite{danvy1992pragmatics}).
The captured variables become formal parameters of $\hat{c}$; their
trust levels are inherited from the judgments at their coordinates via
the transport morphism.  The precondition conjunction ensures both
the body contract and the captured-variable trust requirements are met.
\end{proof}

\subsection{Example: Memoized Closure}

\begin{example}\label{ex:memoize}
Consider the pattern:
\begin{lstlisting}[style=jugeo-python]
cache = {}          # trust: solver_discharged

def memoize(f):
    def wrapper(*args):
        if args not in cache:
            cache[args] = f(*args)
        return cache[args]
    return wrapper
\end{lstlisting}
The closure \texttt{wrapper} captures \texttt{cache} and \texttt{f}.
Closure analysis extracts two captured coordinates:
$\capture(\texttt{wrapper}) = \{\coord_{\texttt{cache}},\,
\coord_{\texttt{f}}\}$.
The contract of \texttt{wrapper} requires
$\tau(\coord_{\texttt{cache}}) \ge \Tsolver$ and
$\tau(\coord_{\texttt{f}}) \ge \Tsolver$, matching the existing
judgment on \texttt{cache}.
\end{example}

% =====================================================================
\section{Decorator Verification}\label{sec:decorator}
% =====================================================================

\Python{} decorators are syntactic sugar for higher-order function
application.  The declaration:
\begin{lstlisting}[style=jugeo-python]
@decorator
def f(x):
    ...
\end{lstlisting}
is exactly equivalent to \texttt{f = decorator(f)}.  From the
geometric perspective, a decorator is a \emph{morphism transformer}:
it maps the morphism $\phi_f$ corresponding to \texttt{f} to a
new morphism $\Decor(\phi_f)$ that routes control through the
decorator wrapper.

\subsection{Decorators as Morphism Composition}

\begin{definition}[Decorator Morphism]\label{def:decormor}
Let $\phi_f : C_{\mathrm{call}} \to C_f$ be the site morphism
for function \texttt{f}, and let
$\phi_d : C_f \to C_{\Decor(f)}$ be the morphism introduced by the
decorator \texttt{d}.  The \emph{decorated morphism} is the
composition
\[
  \phi_{\Decor(f)} = \phi_d \compose \phi_f :
  C_{\mathrm{call}} \to C_{\Decor(f)}.
\]
\end{definition}

The trust level of the decorated morphism is the meet of the
constituent trust levels:

\begin{proposition}[Decorator Trust Propagation]\label{prop:decortrust}
$\tau(\phi_{\Decor(f)}) = \tau(\phi_d) \tatten \tau(\phi_f)$
where $\tatten$ is the trust attenuation operator.
\end{proposition}

\begin{proof}
Trust attenuation is defined as the conservative meet in the trust
lattice $\Talg$ (Paper~04, \S3.2).  Since
$\phi_{\Decor(f)} = \phi_d \compose \phi_f$ routes evidence through
both morphisms, the resulting trust cannot exceed the minimum of the
two; $\tatten$ is defined as exactly this minimum.
\end{proof}

\subsection{Contract Derivation for Decorated Functions}

\begin{construction}[Decorator Contract Rule]\label{const:decorcon}
Given contracts $\Contract(\phi_f) = (\Pre_f, \Post_f, \tau_f)$ and
$\Contract(\phi_d) = (\Pre_d, \Post_d, \tau_d)$, the contract of the
decorated function is:
\[
  \Contract(\phi_{\Decor(f)}) =
    \Bigl(\Pre_d[\hat{f} \mapsto \phi_f] \wedge \Pre_f,\;
          \Post_d[\hat{f} \mapsto \phi_f] \wedge \Post_f,\;
          \tau_d \tatten \tau_f\Bigr)
\]
where the substitution $[\hat{f} \mapsto \phi_f]$ replaces the
abstract function parameter of the decorator with the concrete
morphism for \texttt{f}.
\end{construction}

\subsection{Stacked Decorators}

When multiple decorators are stacked, they compose left-to-right:
\[
  \phi_{\Decor_n \circ \cdots \circ \Decor_1(f)} =
  \phi_{d_n} \compose \cdots \compose \phi_{d_1} \compose \phi_f.
\]
Trust propagates by iterated attenuation:
$\tau = \tau_{d_n} \tatten \cdots \tatten \tau_{d_1} \tatten \tau_f$.

\begin{example}\label{ex:stackeddec}
\begin{lstlisting}[style=jugeo-python]
@require_auth          # trust: solver_discharged
@rate_limit(100)       # trust: runtime_witnessed
def api_endpoint(req):
    ...
\end{lstlisting}
The site contains morphisms
$\phi_{\texttt{require\_auth}} \compose \phi_{\texttt{rate\_limit}} \compose
\phi_{\texttt{api\_endpoint}}$.
Trust is attenuated to $\Truntime$ (the minimum of the three
levels), reflecting that the rate-limit decorator is witnessed only
at run time.
\end{example}

% =====================================================================
\section{Callback Chains}\label{sec:callbacks}
% =====================================================================

Event-driven Python code (Flask routes, asyncio handlers, tkinter
callbacks) organises control flow as \emph{callback chains}: sequences
of callable surfaces through which a single event propagates.  Verifying
such code requires reasoning about the whole chain, not individual links.

\subsection{Sequential Morphism Model}

\begin{definition}[Callback Chain]\label{def:cbchain}
A \emph{callback chain} of length $n$ is a sequence of callable
surfaces $\Chain = (\surf_1, \surf_2, \dots, \surf_n)$ together with
a set of \emph{link conditions}
$L_i : \Post(\surf_i) \to \Pre(\surf_{i+1})$, $1 \le i < n$, stating
that the postcondition of each step implies the precondition of the
next.
\end{definition}

The chain gives rise to a sequence of site morphisms:
\[
  C_0 \xrightarrow{\phi_{\surf_1}} C_1
      \xrightarrow{\phi_{\surf_2}} C_2
      \cdots
      \xrightarrow{\phi_{\surf_n}} C_n.
\]

\begin{proposition}[Chain Composability]\label{prop:chaincomp}
A callback chain $\Chain$ is \emph{composable} if and only if the
link conditions $L_1, \dots, L_{n-1}$ are all satisfied.  In that
case, the chain is equivalent to a single morphism
$\phi_\Chain = \phi_{\surf_n} \compose \cdots \compose \phi_{\surf_1}$
with contract
$\Contract(\phi_\Chain) = (\Pre(\surf_1), \Post(\surf_n),
\tau_{\surf_1} \tatten \cdots \tatten \tau_{\surf_n})$.
\end{proposition}

\begin{proof}
By induction on $n$.  Base case $n=1$ is trivial.  For the inductive
step, the link condition $L_{n-1}$ guarantees that the composition
$\phi_{\surf_n} \compose \phi_{\surf_{n-1} \compose \cdots}$ is
well-defined as a site morphism; trust attenuation is associative in
$\Talg$ (Paper~04, Lemma~3.4).
\end{proof}

\subsection{Asynchronous Callbacks}

\Python's \texttt{async}/\texttt{await} syntax introduces
\emph{suspended morphisms}: site morphisms whose execution may be
interleaved with other event-loop callbacks.

\begin{definition}[Suspended Morphism]\label{def:suspended}
A \emph{suspended morphism} $\phi^\dagger$ is a morphism in the
semantic site equipped with a \emph{suspension point}
$s \in C_{\mathrm{mid}}$ such that
$\phi^\dagger = \phi_{\mathrm{resume}} \compose \phi_{\mathrm{pre}}$
with $\phi_{\mathrm{pre}} : C_0 \to C_{\mathrm{mid}}$ and
$\phi_{\mathrm{resume}} : C_{\mathrm{mid}} \to C_n$.
\end{definition}

Contract verification for suspended morphisms follows the same
attenuation rule, with the additional requirement that the
intermediate state $C_{\mathrm{mid}}$ carries a judgment of trust
$\ge \Truntime$ so that runtime witnesses can validate the suspension
and resumption.

\subsection{Promise-Style Chains}

Promise-style code (e.g.\ via \texttt{asyncio.gather} or
\texttt{concurrent.futures}) creates \emph{parallel callback chains}
whose results are joined at a synchronisation point.

\begin{construction}[Parallel Chain Composition]\label{const:parallel}
Given parallel chains $\Chain_1, \Chain_2$ with a join
function $j : C_{n_1} \times C_{n_2} \to C_{\mathrm{out}}$,
the combined chain has contract:
\[
  \Contract(\Chain_1 \parallel \Chain_2) =
    \Bigl(\Pre(\Chain_1) \wedge \Pre(\Chain_2),\;
          \Post(j),\;
          \tau(\Chain_1) \tatten \tau(\Chain_2) \tatten \tau(j)\Bigr).
\]
\end{construction}

% =====================================================================
\section{Site Construction}\label{sec:site}
% =====================================================================

We now describe how callable surfaces are assembled into a semantic
site, making the construction from \S\ref{sec:model}--\S\ref{sec:callbacks}
precise.

\subsection{The Callable Surface Category}

\begin{definition}[Callable Surface Category]\label{def:surfcat}
The \emph{callable surface category} $\mathbf{Surf}$ has:
\begin{itemize}
  \item \textbf{Objects}: callable surfaces $\surf \in \CS$.
  \item \textbf{Morphisms}: for each pair $(\surf, \surf')$ where
        $\surf$ may directly invoke $\surf'$ (i.e.\ the
        postcondition of $\surf$ implies the precondition of
        $\surf'$), a morphism $\surf \to \surf'$.
  \item \textbf{Composition}: sequential invocation composition,
        with trust propagation by attenuation.
  \item \textbf{Identity}: the identity morphism on $\surf$ is the
        trivial invocation that leaves all contracts unchanged.
\end{itemize}
\end{definition}

\subsection{The Embedding Functor}

\begin{construction}[Embedding Functor]\label{const:functor}
The \emph{callable-surface embedding} is a functor
$\ContEnv : \mathbf{Surf} \to \Site$ defined as follows.
On objects, $\ContEnv(\surf) = C_\surf$ is the coordinate in
$\Site$ corresponding to the function definition containing $\surf$.
On morphisms $f : \surf \to \surf'$, $\ContEnv(f) = \SiteMor(\surf, \surf')$
is the site morphism whose trust is $\tau(\surf) \tatten \tau(\surf')$.
\end{construction}

\begin{proposition}[Functor Laws]\label{prop:functorlaw}
$\ContEnv$ preserves identities and composition:
\begin{enumerate}
  \item $\ContEnv(\mathrm{id}_\surf) = \mathrm{id}_{C_\surf}$.
  \item $\ContEnv(g \compose f) = \ContEnv(g) \compose \ContEnv(f)$.
\end{enumerate}
\end{proposition}

\begin{proof}
(1) The identity morphism on $\surf$ maps to the identity section on
$C_\surf$, which is the identity in the site.
(2) Composition of site morphisms corresponds to sequential
invocation; trust attenuation distributes over composition in $\Talg$
because $\tatten$ is associative and monotone
(Paper~04, \S3.2).
\end{proof}

\subsection{Covering Families from Callable Surfaces}

A set of callable surfaces forms a \emph{covering family} in the
Grothendieck topology of the semantic site when their collective
contracts cover the full precondition space of the surrounding
function.

\begin{definition}[Surface Cover]\label{def:surfcover}
Let $F$ be a Python function with coordinate $C_F$.  A set
$\{\surf_1, \dots, \surf_k\} \subseteq \CS$ of callable surfaces
inside $F$ is a \emph{surface cover} if:
$\Pre(\surf_1) \vee \cdots \vee \Pre(\surf_k) \models \Pre(F)$.
\end{definition}

\begin{proposition}[Coverage Descent]\label{prop:covdescent}
If $\{\surf_1,\dots,\surf_k\}$ is a surface cover for $F$ and each
$\surf_i$ is contractually sound, then $F$ is contractually sound.
\end{proposition}

\begin{proof}
By the sheaf condition on the semantic site (Paper~01, Theorem~4.7):
local sections (contracts at callable surfaces) that agree on
overlaps glue to a global section (contract at $F$).  The surface
cover provides the covering family; contractual soundness at each
surface provides the local sections; the descent check verifies
agreement on overlaps.
\end{proof}

% =====================================================================
\section{Higher-Order Soundness Theorem}\label{sec:theorem}
% =====================================================================

We now state and prove the main theoretical result of the paper.

\subsection{Setup}

Let $P$ be a Core Python program and let $\CS(P)$ be its callable
surface set produced by $\textsc{ExtractSurfaces}$.  A
\emph{complete contract environment} for $P$ is a mapping
$\CEnv : \CS(P) \to \Contract$ such that every surface is
contractually sound:
$\forall \surf \in \CS(P).\; \sound(\surf, \CEnv)$.

An \emph{escape path} is a finite sequence of callable surfaces
$(\surf_1, \surf_2, \dots, \surf_m)$ such that $\surf_1$ is a
HOF argument surface, $\surf_m$ is a direct call surface, and
the chain $(\surf_1, \dots, \surf_m)$ is composable---but some
$\surf_i$ has no contract in $\CEnv$.  Such a path represents a
``contract escape'': a callable value that reaches an unchecked
invocation without passing through any verified morphism.

\subsection{Main Theorem}

\begin{theorem}[Higher-Order Soundness]\label{thm:hosound}
Let $P$ be a Core Python program and let $\CEnv$ be a complete
contract environment for $P$.  Then $P$ has no escape paths: for
every reachable sequence of callable surfaces
$(\surf_1, \surf_2, \dots, \surf_m)$, every surface $\surf_i$ is
contractually sound.
\end{theorem}

\begin{proof}
We prove by contradiction.  Suppose an escape path
$(\surf_1, \dots, \surf_m)$ exists with $\surf_j$ having no
contract in $\CEnv$ for some $1 \le j \le m$.  By
Proposition~\ref{prop:extract}, $\surf_j \in \CS(P)$.  But $\CEnv$
is complete, so every surface in $\CS(P)$ is contractually sound, a
contradiction.

More constructively: the reachable surfaces of any HOF application
are exactly $\{\texttt{hof}\} \cup \texttt{args} \cup
\{\texttt{result}\}$ (Definition~\ref{def:surface}).  By
Proposition~\ref{prop:extract}, all these surfaces appear in
$\CS(P)$.  Since $\CEnv$ is complete, all are contractually sound.
Composability follows from the link conditions inherited from
Proposition~\ref{prop:chaincomp}, and the composed morphism in
$\Site$ is well-defined by the functor laws
(Proposition~\ref{prop:functorlaw}).  Therefore no unchecked callable
can be reached.
\end{proof}

\begin{corollary}[Callback Escape Freedom]\label{cor:cbescape}
Under the same hypotheses as Theorem~\ref{thm:hosound}, every
callback chain in $P$ is composable, and no callback invocation
bypasses a verified contract.
\end{corollary}

\begin{proof}
Callback surfaces are a subclass of callable surfaces.  Completeness
of $\CEnv$ covers them; composability follows from chain link
conditions.
\end{proof}

\begin{corollary}[Decorator Safety]\label{cor:decsafety}
Under the same hypotheses, every decorator application preserves the
trust level of the decorated function up to attenuation.  No decorator
can silently drop a verified contract.
\end{corollary}

\begin{proof}
Decorator morphisms are constructed by Proposition~\ref{prop:decortrust}
to propagate trust through $\tatten$.  Since $\tatten$ is monotone and
the original contract is verified (trust $\ge \Tsolver$), the
decorated trust is $\tau_d \tatten \tau_f \ge \Truntime$ as long as
the decorator itself has trust $\ge \Truntime$.  No silent promotion
is possible (Common.lean, \S5).
\end{proof}

\subsection{Lean Mechanisation}

The proof is mechanised in \LEAN{} in the file
\texttt{proofs/lean/Paper45\_CallableSurfaces.lean} within the
\texttt{JudgmentGeometry.CallableSurfaces} namespace.  The
formalisation defines:
\begin{itemize}
  \item \texttt{SurfaceKind}, \texttt{CallableSurface},
        \texttt{Contract}: the core data types.
  \item \texttt{ContractEnv}: a finite map from surface ids to
        contracts.
  \item \texttt{contractuallySound}: the predicate
        $\sound(\surf, \CEnv)$.
  \item \texttt{HOFApplication}, \texttt{reachableSurfaces}: the
        higher-order application record and its reachable surface set.
  \item \texttt{higher\_order\_soundness}: the main theorem,
        proved without \texttt{sorry}.
\end{itemize}

% =====================================================================
\section{Experiments}\label{sec:experiments}
% =====================================================================

\subsection{Benchmark}

We evaluate \CallSurf{} on a benchmark of $\numcases{}$ \Python{}
programs drawn from $30$ open-source projects including
Flask~\cite{flask}, asyncio standard library~\cite{asyncio}, and the
\texttt{requests} library~\cite{requests}.  The benchmark is stratified
by pattern type:
\begin{itemize}
  \item $\numhof{}$ programs containing higher-order function patterns
        (\MapF, \FilterF, \ReduceF, \texttt{sorted} with key, etc.).
  \item $\numcb{}$ programs containing callback chains (Flask routes,
        asyncio coroutines, \texttt{threading.Thread} targets).
  \item $300 - \numhof{} - \numcb{} = 159$ programs containing
        decorator applications.
\end{itemize}

\subsection{Metrics}

We measure:
\begin{enumerate}
  \item \textbf{Recall}: fraction of callable surfaces correctly
        identified by \textsc{ExtractSurfaces}.
  \item \textbf{Contract coverage}: fraction of identified surfaces
        for which a contract is derived automatically.
  \item \textbf{Soundness rate}: fraction of programs for which
        Theorem~\ref{thm:hosound} applies (complete $\CEnv$).
  \item \textbf{Latency}: wall-clock time per program.
\end{enumerate}

\subsection{Results}

\begin{table}[H]
\centering
\caption{Callable Surface Analysis results by pattern category.
         HOF = higher-order function; CB = callback chain;
         Dec = decorator.}\label{tab:results}
\begin{tabular}{lrrrr}
\toprule
\textbf{Category} & \textbf{Programs} & \textbf{Recall (\%)}
  & \textbf{Coverage (\%)} & \textbf{Latency (ms)} \\
\midrule
HOF          & \numhof   & \hofrecall  & \ppFortyFiveHofCoverage & \medtime  \\
Callback (CB) & \numcb    & \cbrecall   & \ppFortyFiveCbCoverage & \medtime  \\
Decorator    & 159        & \ppFortyFiveDecRecall         & \ppFortyFiveDecCoverage & \medtime  \\
\midrule
\textbf{Overall} & \numcases & \multicolumn{2}{c}{\emph{See text}} & \textbf{\medtime} \\
\bottomrule
\end{tabular}
\end{table}

Callable surface recall is high across all categories.
Decorators achieve the highest recall because their syntactic
form is highly regular; callbacks achieve the lowest due to
dynamic registration patterns in event-loop code that are not yet
handled by the current callee classifier.

Contract coverage is high overall.  The gap consists
primarily of dynamically constructed callables (functions created by
\texttt{type()} or \texttt{functools.partial}) for which the analyzer
falls back to a conservative unverified contract.  Programs with full
contract coverage all satisfy the Higher-Order Soundness Theorem
(proved in \LEAN{}); the remainder have at least one surface with a
conservative contract below $\Tsolver$.

\subsection{Comparison with Baselines}

We compare against three baselines:
\begin{itemize}
  \item \textbf{mypy}~\cite{lehtosalo2016mypy}: type-based analysis;
        does not verify contracts, only types.  Reports $0\%$ contract
        coverage by definition.
  \item \textbf{Pyright}~\cite{pyright}: similar to mypy.
  \item \textbf{Pyre}~\cite{pyre}: Meta's type checker; richer type
        inference but still no contract verification.
\end{itemize}

All three baseline tools correctly identify callable surfaces (via
call-graph construction) but do not attach contracts to them.
\CallSurf{} is the only system that (a) maps surfaces to site morphisms,
(b) derives contracts from the judgment algebra, and (c) proves
higher-order soundness via the completeness theorem.

\subsection{Ablation Study}

Disabling closure analysis reduces contract coverage
on programs with closures.  Disabling decorator
verification degrades soundness rate.  Disabling
callback chain composition reduces recall on asyncio programs.  These ablations confirm that all three sub-analyses
contribute meaningfully to overall performance.

% =====================================================================
\section{Conclusion}\label{sec:conclusion}
% =====================================================================

We presented \emph{callable surface analysis}, a framework for
verifying higher-order \Python{} functions within the \JuGeo{}
semantic site model.  The framework identifies callable surfaces
in the AST, maps them to site morphisms, and derives contracts
via three specialised analyses for closures, decorators, and
callback chains.

The central contribution is the Higher-Order Soundness Theorem
(Theorem~\ref{thm:hosound}): under a complete contract environment,
no callable value can escape without passing through a verified
morphism.  The theorem is proved formally in \LEAN{} without
\texttt{sorry}, providing a machine-checked guarantee that complements
the empirical evaluation.

Experiments on $\numcases{}$ programs confirm high surface recall
and contract coverage at a median latency of $\medtime{}$
per program, consistent with the universal benchmark (\expAccuracy{}
accuracy on \expTotalPrograms{} programs).

\paragraph{Limitations.}
The current callee classifier does not handle all forms of dynamic
callable construction; some surfaces receive conservative
contracts.  The link-condition checker for callback chains requires
explicit postcondition annotations, which are not always present.

\paragraph{Future work.}
Natural extensions include: (a)~a \emph{lazy callable surface
refinement} pass that upgrades conservative contracts when runtime
witnesses are available; (b)~integration with the Treaty Synthesis
framework (Paper~08) to synthesise missing contracts from usage
patterns; (c)~extension to \texttt{eval}-containing programs via
a sandboxed execution environment.

\paragraph{Acknowledgements.}
The author thanks the \JuGeo{} development community for benchmark
contributions and for feedback on early drafts.

\begin{thebibliography}{99}

\bibitem{jg01}
H.~Young.
\newblock Grothendieck topologies for compositional program analysis.
\newblock \textit{Judgment Geometry Series}, Paper~01, 2024.

\bibitem{jg04}
H.~Young.
\newblock Trust algebras for compositional verification.
\newblock \textit{Judgment Geometry Series}, Paper~04, 2024.

\bibitem{jg07}
H.~Young.
\newblock Verifying effectful Python without leaving Python.
\newblock \textit{Judgment Geometry Series}, Paper~07, 2024.

\bibitem{jg08}
H.~Young.
\newblock Treaty synthesis for cross-module verification.
\newblock \textit{Judgment Geometry Series}, Paper~08, 2024.

\bibitem{jg38}
H.~Young.
\newblock Semantic caching and incremental verification.
\newblock \textit{Judgment Geometry Series}, Paper~38, 2024.

\bibitem{danvy1992pragmatics}
O.~Danvy and A.~Filinski.
\newblock Representing control: A study of the CPS transformation.
\newblock \textit{Mathematical Structures in Computer Science},
  2(4):361--391, 1992.

\bibitem{church1936unsolvable}
A.~Church.
\newblock An unsolvable problem of elementary number theory.
\newblock \textit{American Journal of Mathematics}, 58(2):345--363, 1936.

\bibitem{reynolds1972definitional}
J.~C.~Reynolds.
\newblock Definitional interpreters for higher-order programming languages.
\newblock In \textit{Proceedings of ACM Annual Conference}, pages 717--740,
  1972.

\bibitem{mitchell1988abstract}
J.~C.~Mitchell and G.~D.~Plotkin.
\newblock Abstract types have existential type.
\newblock \textit{ACM Transactions on Programming Languages and Systems},
  10(3):470--502, 1988.

\bibitem{milner1978theory}
R.~Milner.
\newblock A theory of type polymorphism in programming.
\newblock \textit{Journal of Computer and System Sciences},
  17(3):348--375, 1978.

\bibitem{flanagan2002extended}
C.~Flanagan, K.~R.~M.~Leino, M.~Lillibridge, G.~Nelson, J.~B.~Saxe,
  and R.~Stata.
\newblock Extended static checking for Java.
\newblock In \textit{PLDI 2002}, pages 234--245, 2002.

\bibitem{lehtosalo2016mypy}
J.~Lehtosalo et~al.
\newblock Mypy: Optional static typing for Python.
\newblock \url{https://mypy-lang.org}, 2016.

\bibitem{pyright}
Microsoft.
\newblock Pyright: Static type checker for Python.
\newblock \url{https://github.com/microsoft/pyright}, 2019.

\bibitem{pyre}
Meta Platforms.
\newblock Pyre: A performant type-checker for Python~3.
\newblock \url{https://pyre-check.org}, 2018.

\bibitem{flask}
A.~Ronacher.
\newblock Flask: A micro web framework for Python.
\newblock \url{https://flask.palletsprojects.com}, 2010.

\bibitem{asyncio}
G.~van~Rossum.
\newblock \texttt{asyncio} --- Asynchronous I/O.
\newblock \textit{Python Standard Library Documentation}, 2014.

\bibitem{requests}
K.~Reitz.
\newblock Requests: HTTP for Humans.
\newblock \url{https://requests.readthedocs.io}, 2011.

\bibitem{moura2021lean4}
L.~de~Moura and S.~Ullrich.
\newblock The Lean 4 theorem prover and programming language.
\newblock In \textit{CADE-28}, LNCS 12699, pages 625--635, 2021.

\bibitem{grothendieck1972sga}
A.~Grothendieck et~al.
\newblock \textit{Th\'eorie des Topos et Cohomologie \'Etale des
  Sch\'emas}, Lecture Notes in Mathematics 269/270/305.
\newblock Springer, 1972.

\bibitem{johnstone2002sketches}
P.~T.~Johnstone.
\newblock \textit{Sketches of an Elephant: A Topos Theory Compendium}.
\newblock Oxford University Press, 2002.

\end{thebibliography}

\end{document}
